\titleframe

% Contoh Kode Slide
%\section{Introduction}
%\begin{frame}{\insertsectionhead}
%	\framesubtitle{A short introduction to Trigon}
%	\themename is a modern, elegant and versatile theme for Beamer, inspired by
%	the
%	\href{https://github.com/matze/mtheme}{\textsc{metropolis} theme} from Matthias
%	Vogelgesang.
%	\vfill
%	\themename comes with lots of nice extra features
%	\begin{itemize}
%		\item Multiple style variations for title, section and normal slides
%		\item Simple customization of theme colors
%		\item Lots of convenient options to tweak the design
%	\end{itemize}
%\end{frame}

% Contoh Kode Tabel
%  \begin{table}[H]
% \centering
% \caption{A nice table example}
% \begin{tabular}{@{} lccc @{}}
%    \toprule
%    & \textbf{Velocity} & \textbf{Angle}  & \textbf{Vertical force} \\
%    & $U$ & $\alpha$  & $F_z$ \\
%    & [m/s] & [$^\circ$]  & [N] \\
%    \midrule
%    2D simulation  & 9 & 2 & 9.23 \\
%    3D simulation  & 10.0 & 3 & 15.039 \\
%    Experiment A   & 11.31 & 2.5 & 13.2 \\
%    Experiment B   & 11.26 & 2.7 & 12.6 \\
%    Experiment C   & 11.33 & 2.47 & 13.6 \\
%    \bottomrule
%  \end{tabular}
%\end{table}

% Contoh Kode Gambar
%  \begin{figure}[ht!]
%	\begin{subfigure}[b]{0.3\textwidth}
%		\frame{\includegraphics[width=\textwidth]{layout_example-03.jpg}}
%		\caption*{plain}
%	\end{subfigure}
%	\hspace{\fill}
%	\begin{subfigure}[b]{0.3\textwidth}
%		\frame{\includegraphics[width=\textwidth]{layout_example-02.jpg}}
%		\caption*{style1}
%	\end{subfigure}
%	\hspace{\fill}
%	\begin{subfigure}[b]{0.3\textwidth}
%		\frame{\includegraphics[width=\textwidth]{layout_example-01.jpg}}
%		\caption*{style2 (default)}
%	\end{subfigure}
%\end{figure}

% Contoh Kode Blocks
% \begin{block}{Regular block}
%	Just a regular block
%\end{block}
%\begin{alertblock}{Alert block}
%	Some important thing
%\end{alertblock}
%\begin{exampleblock}{Example block}
%	No difference with regular block to avoid excessive distraction
%\end{exampleblock}

\setbeamertemplate{frame footer}{Alauddin Maulana Hirzan, M. Kom - Universitas Semarang}

\section{\textit{Internet of Things} dan \textit{Cloud Computing}}
\begin{frame}{\insertsectionhead}
	\framesubtitle{Apa itu \textit{Cloud Computing}? \#1}
	\justifying
	Kita mengenal teknologi \textit{Cloud Computing} atau Komputasi Awan sebagai sebuah teknologi yang mengandalkan platform yang ada di \textbf{Internet} untuk melakukan komputasi maupun penyimpanan data secara permanen.
	\vfill
	\begin{block}{Info}
		Contoh nyata dari platform ini adalah \textit{Google Drive} untuk menyimpan file dan \textit{Google Colab} untuk melakukan analisis dan pengambilan tindakan berdasarkan data di Drive.
	\end{block}
\end{frame}

\begin{frame}{\insertsectionhead}
	\framesubtitle{Apa itu \textit{Cloud Computing}? \#2}
	\justifying
	Menurut \textbf{ibm.com}, \textit{Cloud Computing} adalah
	\vfill
	\begin{quote}
		Akses berdasarkan permintaan, melalui internet, ke sumber daya komputasi—aplikasi, server (server fisik dan server virtual), penyimpanan data, alat pengembangan, kemampuan jaringan, dan banyak lagi—dihosting di pusat data jarak jauh yang dikelola oleh layanan penyedia cloud (atau CSP).
	\end{quote}
\end{frame}

\begin{frame}{\insertsectionhead}
	\framesubtitle{Apa itu \textit{Cloud Computing}? \#3}
	\justifying
	Dibandingkan dengan TI lokal tradisional, dan bergantung pada layanan cloud yang Anda pilih, komputasi awan membantu melakukan hal berikut:
	\vfill
	\begin{itemize}
		\item Biaya TI yang lebih rendah
		\item Tingkatkan kelincahan dan nilai waktu
		\item Skala lebih mudah dan hemat biaya
	\end{itemize}
\end{frame}

\begin{frame}{\insertsectionhead}
	\framesubtitle{Apa itu \textit{Cloud Computing}? \#4 - Biaya Rendah}
	\justifying
	Cloud memungkinkan pengguna melepas sebagian atau sebagian besar biaya dan upaya untuk membeli, menginstal, mengonfigurasi, dan mengelola infrastruktur lokal pengguna itu sendiri.
	\vfill
	Sehingga pengguna dapat memiliki layanan yang dapat diakses di mana saja selama ada Internet di sana.
	\vfill
	\begin{block}{Info}
		\textbf{Google Drive} merupakan contoh nyata pengganti flash disk yang gratis dimiliki siapa saja
	\end{block}
\end{frame}

\begin{frame}{\insertsectionhead}
	\framesubtitle{Apa itu \textit{Cloud Computing}? \#5 - Agilitas Layanan}
	\justifying
	Dengan cloud, organisasi pengguna dapat mulai menggunakan aplikasi perusahaan dalam hitungan menit, alih-alih menunggu berminggu-minggu atau berbulan-bulan agar TI merespons permintaan, membeli dan mengonfigurasi perangkat keras pendukung, serta menginstal perangkat lunak.
	\vfill
	\begin{block}{Info}
		TInggi tingkat respon dari \textit{Cloud Computing} inilah yang juga menjadi daya tarik pengguna level industri
	\end{block}
\end{frame}

\begin{frame}{\insertsectionhead}
	\framesubtitle{Apa itu \textit{Cloud Computing}? \#6 - Peningkatan Perangkat}
	\justifying
	Cloud memberikan elastisitas—daripada membeli kelebihan kapasitas yang tidak terpakai selama periode lambat, Pengguna dapat meningkatkan dan menurunkan kapasitas sebagai respons terhadap lonjakan dan penurunan lalu lintas.
	\vfill
	\begin{block}{Info}
		Fitur ini membantu menekan pengeluaran yang diperlukan oleh industri dalam menghadirkan suatu layanan.
	\end{block}
\end{frame}

\begin{frame}{\insertsectionhead}
	\framesubtitle{Apa itu \textit{Cloud Computing}? \#7}
	\justifying
	Istilah 'komputasi awan' juga mengacu pada teknologi yang membuat cloud bekerja. Ini mencakup beberapa bentuk infrastruktur TI tervirtualisasi—server, perangkat lunak sistem operasi, jaringan, dan infrastruktur lain yang diabstraksikan, menggunakan perangkat lunak khusus, sehingga dapat dikumpulkan dan dibagi terlepas dari batasan perangkat keras fisik.
	\vfill
	Misalnya, server perangkat keras tunggal dapat dibagi menjadi beberapa server virtual.
	\vfill
	\begin{block}{Info}
		Teknologi inilah yang disebut dengan teknologi virtualisasi
	\end{block}
\end{frame}

\begin{frame}{\insertsectionhead}
	\framesubtitle{Layanan \textit{Cloud Computing}? \#1}
	\justifying
	Layanan yang disediakan oleh Cloud dapat berupa:
	\vfill
	\begin{itemize}
		\item SaaS (Software-as-a-Service)
		\item PaaS (Platform-as-a-Service)
		\item IaaS (Infrastructure-as-a-Service)
	\end{itemize}
	\vfill
	Layanan yang digunakan bergantung dari yang ingin dimanfaatkan.
\end{frame}

\begin{frame}{\insertsectionhead}
	\framesubtitle{Layanan \textit{Cloud Computing}? \#2 - \textit{Software-as-a-Service}}
	\justifying 
	Layanan ini adalah perangkat lunak aplikasi yang dihosting di cloud dan pengguna gunakan melalui browser web, klien desktop khusus, atau API yang terintegrasi dengan sistem operasi desktop atau seluler. 
	\vfill
	Dalam kebanyakan kasus, pengguna SaaS membayar biaya berlangganan bulanan atau tahunan; beberapa mungkin menawarkan harga 'bayar sesuai pemakaian' berdasarkan penggunaan yang sebenarnya.
	\vfill
	\begin{block}{Info}
		Google Doc merupakan salah satu contoh aplikasi sebagai service yang bisa digunakan oleh pengguna melalui browser atau seluler
	\end{block}
\end{frame}

\begin{frame}{\insertsectionhead}
	\framesubtitle{Layanan \textit{Cloud Computing}? \#2 - \textit{Platform-as-a-Service}}
	\justifying 
	Layanan ini menyediakan pengembang perangkat lunak dengan platform sesuai permintaan—perangkat keras, tumpukan perangkat lunak lengkap, infrastruktur, dan bahkan alat pengembangan—untuk menjalankan, mengembangkan, dan mengelola aplikasi tanpa biaya, kerumitan, dan ketidakfleksibelan pemeliharaan platform tersebut di tempat.
	\vfill
	\begin{block}{Info}
		Blynk, Cayenne, Antares merupakan platform yang disediakan untuk Plaform \textit{Internet of Things}
	\end{block}
\end{frame}

\begin{frame}{\insertsectionhead}
	\framesubtitle{Layanan \textit{Cloud Computing}? \#3 - \textit{Infrastructure-as-a-Service}}
	\justifying 
	IaaS menyediakan akses sesuai permintaan ke sumber daya komputasi dasar-server fisik dan virtual, jaringan, dan penyimpanan-melalui internet dengan basis bayar sesuai pemakaian.
	\vfill
	\begin{block}{Info}
		Virtual Private Server merupakan contoh dari IaaS yang menyediakan keseluruhan infrastruktur sebagai layanan
	\end{block}
\end{frame}

\begin{frame}{\insertsectionhead}
	\framesubtitle{\textit{Cloud Computing} dan \textit{Internet of Things} \#1}
	\justifying 
	Salah satu komponen yang meningkatkan keberhasilan Internet of Things adalah Cloud Computing. Komputasi awan memungkinkan pengguna untuk melakukan tugas komputasi menggunakan layanan yang disediakan melalui Internet. 
	\vfill
	Penggunaan Internet of Things dalam hubungannya dengan teknologi cloud telah menjadi semacam katalis: Internet of Things dan komputasi awan sekarang saling terkait. Ini adalah teknologi masa depan sejati yang akan membawa banyak manfaat.
\end{frame}

\begin{frame}{\insertsectionhead}
	\framesubtitle{\textit{Cloud Computing} dan \textit{Internet of Things} \#2}
	\justifying 
	Karena pesatnya pertumbuhan teknologi, masalah penyimpanan, pemrosesan, dan pengaksesan data dalam jumlah besar telah muncul. Inovasi hebat berkaitan dengan saling menggunakan Internet of Things dan teknologi cloud. Dalam kombinasi, dimungkinkan untuk menggunakan pemrosesan yang kuat dari aliran data sensorik dan layanan pemantauan baru
\end{frame}

\begin{frame}{\insertsectionhead}
	\framesubtitle{\textit{Cloud Computing} dan \textit{Internet of Things} \#3}
	\justifying 
	Hubungan apa yang dipunyai oleh \textit{Cloud Computing} berkaitan dengan \textit{Internet of Things}
	\vfill
	\begin{itemize}
		\item Penyimpanan Data karena terbatasnya perangkat IoT
		\item Keamanan Data yang harus dilindungi dari pihak luar
	\end{itemize}
	\vfill
	Karena alasan tersebut menjadi faktor menentukan untuk menggunakan teknologi \textit{Cloud Computing} dalam perangkat \textit{Internet of Things}
\end{frame}

\begin{frame}{\insertsectionhead}
	\framesubtitle{\textit{Cloud Computing} dan \textit{Internet of Things} \#4 - \textit{Data Storage} \#1}
	\justifying 
	Untuk memahami peran komputasi awan di IoT, mari kita pikirkan situasi di mana kita mengeluarkan beberapa komponen dari sistem seperti tempat penyimpanan dan menyimpannya di awan. Apakah hal ini akan menggangu kinerja sistem? Sistem tidak akan terganggu bahkan sistem memiliki tempat alternatif untuk menyimpan data
	\vfill
	Misalnya, Pengguna memprogram Alexa untuk memberi tahu info tentang berita terbaru, cuaca, dan detail lalu lintas setiap pagi. Alexa akan mengirimkan permintaan ke aplikasi cloud seperti Google Maps, dll untuk mengambil informasi dan memberikannya kepada pengguna.
\end{frame}

\begin{frame}{\insertsectionhead}
	\framesubtitle{\textit{Cloud Computing} dan \textit{Internet of Things} \#5 - \textit{Data Storage}}
	\justifying
	  \begin{figure}[ht!]
			\begin{subfigure}[b]{0.7\textwidth}
					\frame{\includegraphics[width=\textwidth]{alexa-pictionary.jpg}}
				\end{subfigure}
		\end{figure}
\end{frame}

\begin{frame}{\insertsectionhead}
	\framesubtitle{\textit{Cloud Computing} dan \textit{Internet of Things} \#6 - \textit{Data Security} \#1}
	\justifying
	integrasi IoT dan Cloud Computing yang aman juga membantu menghindari kebocoran dan serangan data. Menyimpan data (yang mungkin termasuk kunci dan kata sandi) pada perangkat lokal, umumnya bukan ide yang baik. 
	\vfill
	Selain itu, menyimpan sejumlah besar data pada mesin lokal dapat menjadi proses yang membosankan dan mahal. Batas pemrosesan pada mesin lokal juga terbatas, dan untuk operasi dan aplikasi yang kompleks, cloud tampaknya menjadi pilihan yang layak.
	\vfill
	\begin{block}{Info}
		Sangat tidak direkomendasikan untuk menyimpan data sensitif seperti password di perangkat. Sebaiknya data disimpan di \textit{Cloud}
	\end{block}
\end{frame}

\begin{frame}{\insertsectionhead}
	\framesubtitle{\textit{Cloud Computing} dan \textit{Internet of Things} \#7}
	\justifying 
	Manfaat yang dapat dirasakan ketika menggabungkan IoT dengan Cloud:
	\vfill
	\begin{itemize}
		\item   IoT Cloud Computing menyediakan banyak opsi konektivitas, menyiratkan akses jaringan yang besar.
		\item   Pengembang dapat menggunakan komputasi awan IoT sesuai permintaan. 
		\item  Berdasarkan permintaan, pengguna dapat menskalakan layanan sesuai dengan kebutuhan mereka.
		\item  Cloud Computing menyiratkan penyatuan sumber daya yang ada
	\end{itemize}
\end{frame}

\begin{frame}{\insertsectionhead}
	\framesubtitle{Peran \textit{Cloud Computing} untuk \textit{Internet of Things} \#8}
	\justifying 
	\textbf{\textit{Remote Computing Power}}
	\vfill
	Dengan langkah cepat perkembangan 5G dengan kecepatan internetnya, teknologi cloud semakin mudah mengakses layanan komputasi jarak jauh dengan mengklik mouse. Dengan mengurangi kebutuhan untuk memelihara infrastruktur lokal, cloud telah memungkinkan organisasi untuk melampaui aplikasi konvensional IoT
	\vfill
	\begin{block}{Info}
		Semakin kencang internet yang disediakan, memudahkan pengguna untuk mengakses perangkat di manapun dan kapanpun
	\end{block}
\end{frame}

\begin{frame}{\insertsectionhead}
	\framesubtitle{Peran \textit{Cloud Computing} untuk \textit{Internet of Things} \#9}
	\justifying 
	\textbf{\textit{Security and Privacy}}
	\vfill
	Menjamurnya perangkat IoT mungkin memungkinkan organisasi untuk mengotomatisasi tugas, tetapi juga menimbulkan masalah keamanan yang serius. Cloud dengan berbagai kontrolnya dapat menjadi solusi yang layak di sini. 
	\vfill
	Solusi cloud memudahkan penerapan langkah-langkah keamanan yang sangat mudah - memungkinkan perusahaan untuk mengadopsi enkripsi dan protokol otentikasi yang kuat. Dengan solusi cloud terbaik, dimungkinkan untuk mengelola dan mengamankan identitas pengguna yang mengakses perangkat IoT.
\end{frame}

\begin{frame}{\insertsectionhead}
	\framesubtitle{Peran \textit{Cloud Computing} untuk \textit{Internet of Things} \#10}
	\justifying 
	\textbf{\textit{Data Integration}}
	\vfill
	Karena IoT terus menjadi terkenal, perusahaan telah mulai bereksperimen dengan perangkat yang terhubung untuk mengekstrak informasi real-time tentang proses bisnis utama. 
	\vfill
	Sementara perangkat ini meningkatkan efisiensi operasional dan mengoptimalkan biaya, mereka juga menghasilkan data raksasa yang terlalu rumit untuk diproses bahkan untuk platform analitik mereka.
\end{frame}

\begin{frame}{\insertsectionhead}
	\framesubtitle{Peran \textit{Cloud Computing} untuk \textit{Internet of Things} \#11}
	\justifying 
	\textbf{\textit{Low Entry Barrier}}
	\vfill
	Para inovator di domain IoT mencari solusi hosting tanpa kerumitan. Solusi cloud hosting cukup tepat dalam situasi seperti itu. Dengan solusi hosting cloud, pemain IoT memanfaatkan kekuatan pusat data jarak jauh di India tanpa harus menginstal perangkat keras dan perangkat lunak lokal yang rumit. 
	\vfill
	Selain itu, layanan cloud ini beroperasi pada model 'bayar sesuai penggunaan' di mana pengguna ditagih sesuai dengan sumber daya yang dikonsumsinya. Akibatnya, perusahaan dapat menghindari biaya dimuka yang besar.
\end{frame}

\begin{frame}{\insertsectionhead}
	\framesubtitle{Peran \textit{Cloud Computing} untuk \textit{Internet of Things} \#12}
	\justifying 
	\textbf{\textit{Inter-device Communication}}
	\vfill
	Selain berkomunikasi dengan cloud, perangkat dan layanan IoT perlu terhubung satu sama lain. Solusi cloud memfasilitasi komunikasi tanpa batas antara perangkat IoT. Mereka mengaktifkan banyak API yang kuat seperti Cloudflare dan Dropstr dan memungkinkan interaksi antara perangkat yang terhubung dan smartphone sehingga membuka jalan bagi pertumbuhan teknologi yang terhubung.
	\vfill
	\begin{block}{Info}
		Teknologi M-2-M digunakan untuk mengentaskan permasalahan otomatisasi
	\end{block}
\end{frame}

\begin{frame}{\insertsectionhead}
	\framesubtitle{Peran \textit{Cloud Computing} untuk \textit{Internet of Things} \#13}
	\justifying 
	\begin{figure}[ht!]
		\begin{subfigure}[b]{0.7\textwidth}
			\frame{\includegraphics[width=\textwidth]{firebase.jpg}}
		\end{subfigure}
	\end{figure}
\end{frame}

\setbeamertemplate{background} 
{
	\includegraphics[width=\paperwidth,height=\paperheight]{thank-you.jpg}
}
\begin{frame}[plain]
\end{frame}